% ECONOMICS_PROBLEM: {"id":"D63-72","title":"EFM with nonadditive mixed items and desirable cake","jel_code":"D63","source_id":"OP-0594"}
% FIRST_POSTED: 2026-10-11
% ATTEMPT_STATUS: solved
% SUBMISSION: {"kind":"research_attempt","category":"economics","problem_id":"D63-72","model":{"provider":"OpenAI","version":"GPT-6 (Codex)","version_uncertainty":"The Codex session identifies the GPT-6 family; an exact deployed snapshot identifier is not exposed.","reasoning":"Not exposed","generated_on":"2026-10-11"},"other_models":"Mathematical drafting and adversarial checks used Codex subagents with inherited model settings; individual snapshot identifiers were not exposed. Existing catalogue attempts were read for scope, not used as proof dependencies.","tools":"Primary-source web retrieval; standalone LaTeX; standard-library Python exhaustive structural checks on remote Ubuntu/Python 3.12.3; no Lean formalization.","target":"Does every finite normalized doubly monotone item instance with arbitrary nonnegative integrable cake densities admit a complete strong EFM allocation, with item-only deletion fallback?","scope":"Theorem 1 gives six agents and seven strictly monotone nonadditive chores; common cake total C admits strong EFM exactly for C in [0,1] union [2,infinity). C=3/2 disproves the universal original assertion.","human_contributions":"The human supplied the research objective, scope and publication instructions. No independent human mathematical proof verification is documented. GitHub submission by jbaelaw.","verification":"Separate Codex agents audited all allocation cases, the valuation domain, endpoint constructions and source definition. These are model-assisted checks, not independent human verification.","reproducibility":"Final manuscript prepared under repository Prompt A and then audited under Prompt B; public source and verification artifacts in the companion repository. Private conversation transcripts are not published."}
% FULL_SOLUTION_SCOPE: {"target":"Does every finite normalized doubly monotone item instance with arbitrary nonnegative integrable cake densities admit a complete strong EFM allocation, with item-only deletion fallback?","result":"Theorem 1: exact feasible common-cake totals [0,1] union [2,infinity) for the explicit six-agent seven-chore instance; in particular nonexistence at C=3/2.","coverage":"One admissible counterexample negates the universal quantifier. The proof covers every complete item allocation and every measurable cake partition using identical densities; all original domain assumptions are verified. Strong and weak EFM are explicitly distinguished.","dependencies":"No external existence, measure-division or optimization theorem is used. The cited source supplies the definition and original open question; the proof is elementary.","human_verification":"None documented. Independent mathematical checking in this workflow means separate model agents only.","unverified":"No missing step is asserted within the stated hand proof. Independent human review and formal proof-assistant verification have not been performed.","novelty":"Bounded primary-source and repository/PR comparison completed on 2026-10-11; no matching full counterexample found in inspected sources. No exhaustive novelty or priority claim."}
% !TeX program = pdflatex
\documentclass[11pt,a4paper]{article}
\usepackage[margin=26mm]{geometry}
\usepackage[T1]{fontenc}
\usepackage{lmodern,microtype,amsmath,amssymb,amsthm}
\usepackage[hidelinks]{hyperref}
\setlength{\emergencystretch}{2em}
\newtheorem{theorem}{Theorem}
\newtheorem{proposition}[theorem]{Proposition}
\newcommand{\EFM}{\mathrm{EFM}}
\title{Proof of failure of strong EFM\\for nonadditive chores and cake}
\author{OpenAI GPT-6 (Codex)\\\small Submitted by \texttt{@jbaelaw}}
\date{11 October 2026}
\hypersetup{pdftitle={Proof of failure of strong EFM for nonadditive chores and cake},pdfauthor={OpenAI GPT-6 (Codex); submitted by jbaelaw}}
\begin{document}
\maketitle

\begin{abstract}
Strong envy-freeness for mixed resources need not exist for doubly
monotone item valuations. We give six agents, seven chores, and identical
constant cake densities for which no complete allocation satisfies the
strong EFM definition of Aziz, Lu, Mackenzie, and Suzuki. Every item has
strictly negative marginal value. For the same item valuations, we
determine the exact set of feasible cake totals: $[0,1]\cup[2,\infty)$.
The nonexistence proof covers every item allocation and every measurable
cake partition. The distinction between strong and weak EFM is essential.
\end{abstract}

\section{Definition and statement}

Aziz, Lu, Mackenzie, and Suzuki prove existence of EFM allocations for
additive items and desirable cake, and ask whether existence extends to
doubly monotone valuations~\cite[Section~5]{ALMS}.
We use their Definition~2.3, including its item-only deletion condition.
This is the strong EFM convention in the catalogue problem D63-72.

An allocation consists of item bundles $(O_i)_{i\in N}$ partitioning a
finite set $M$, and measurable cake pieces $(D_i)_{i\in N}$ partitioning
$[0,1]$ up to null sets. Agent $i$ has utility
\[
 u_i(O,D)=v_i(O)+V_i(D),\qquad
 V_i(D)=\int_D f_i(t)\,dt.
\]
The original question quantifies over all finite $N,M$, all normalized
item values $v_i(\varnothing)=0$, and all nonnegative Lebesgue-integrable
$f_i$. It assumes that each agent has a fixed partition $M=G_i\sqcup B_i$
such that, for every $S\subseteq M$ and $g\notin S$,
\[
 g\in G_i\ \Longrightarrow\ v_i(S\cup\{g\})\geq v_i(S),\qquad
 g\in B_i\ \Longrightarrow\ v_i(S\cup\{g\})\leq v_i(S).
\]
The sign partition may depend on the agent. Such values are called
\emph{doubly monotone}; either part may be empty.
An allocation is \emph{strong EFM} if, for every ordered pair $i,j$, either
\begin{equation}\label{eq:EF}
 v_i(O_i)+V_i(D_i)\geq v_i(O_j)+V_i(D_j),
\end{equation}
or
\begin{equation}\label{eq:fallback}
 V_i(D_j)=0
 \quad\text{and}\quad
 \exists g\in O_i\cup O_j:\quad
 v_i(O_i\setminus\{g\})\geq v_i(O_j\setminus\{g\}).
\end{equation}
In particular, the last inequality does not include $V_i(D_i)$.

Let
\[
 N=\{X_1,X_2,Y_1,Y_2,Z_1,Z_2\},\qquad
 M=\{x,y,z,d_1,d_2,d_3,d_4\}.
\]
The designated item $a_i$ is $x$ for agents of type $X$, $y$ for type
$Y$, and $z$ for type $Z$. Define $v_i=-c_i$, where, for every $S\subseteq M$,
\begin{equation}\label{eq:cost}
 c_i(S)=
 \begin{cases}
 0,&S=\varnothing,\\
 1,&S=\{a_i\},\\
 2,&|S|=1\text{ and }a_i\notin S,\\
 5,&|S|=2\text{ and }a_i\in S,\\
 3,&|S|=2\text{ and }a_i\notin S,\\
 |S|+3,&|S|\geq3.
 \end{cases}
\end{equation}
All agents have the same constant density $f_i(t)=C$, where $C\geq0$.

\begin{theorem}\label{thm:main}
For the valuations~\eqref{eq:cost}, a complete strong-EFM allocation
exists if and only if
\[
 C\in[0,1]\cup[2,\infty).
\]
Consequently, taking $C=3/2$ gives a counterexample to universal strong-EFM
existence for doubly monotone items and nonnegative integrable cake densities.
\end{theorem}

The instance satisfies the original assumptions without a limiting
argument. Every cost is normalized, and adding any item strictly
increases cost: the increases from sizes $0$ to $1$, $1$ to $2$, and
$2$ to $3$ are at least one; all subsequent increases equal one.
Thus each $v_i$ is strictly decreasing under every item addition.
It is doubly monotone with $G_i=\varnothing$ and $B_i=M$.
The costs are nonadditive: for example,
$c_{X_1}(\{x,y\})=5>c_{X_1}(\{x\})+c_{X_1}(\{y\})=3$.

For any cake partition write
\begin{equation}\label{eq:payments}
 p_j=C\,|D_j|,\qquad p_j\geq0,\qquad \sum_jp_j=C,
\end{equation}
where $|D_j|$ denotes Lebesgue measure. Every agent values $D_j$ at the
same number $p_j$. Conversely, for $C>0$ every nonnegative vector with
sum $C$ is realized by consecutive cake intervals of lengths $p_j/C$.
For $C=0$, every cake partition has value zero.

\section{Nonexistence for \texorpdfstring{$1<C<2$}{1<C<2}}

Suppose a strong-EFM allocation exists, and use~\eqref{eq:payments}.

\paragraph{Every item bundle is nonempty.}
If $O_j=\varnothing$, some other bundle $O_i$ contains at least two of
the seven items. Its cost to its owner is at least three, and hence
\[
 -c_i(O_i)+p_i\leq-3+C<0\leq p_j.
\]
Agent $i$ envies $j$. Deleting any one item from $O_i$ leaves a
nonempty bundle of positive cost, whereas $O_j$ is empty and has no
item to delete. The item-only inequality in~\eqref{eq:fallback}
therefore fails. This excludes every allocation with an empty bundle.

\paragraph{The unique pair contains its owner's designated item.}
Seven items in six nonempty bundles give one pair $P=O_h$ and five
singletons. If $a_h\notin P$, some agent $j$ holds $\{a_h\}$. Then
\[
 -c_h(P)+p_h\leq-3+C<-1
 \leq-c_h(\{a_h\})+p_j.
\]
Deleting either item from $P$ leaves cost two, exceeding the compared
cost one. Deleting the compared item instead leaves compared cost zero
and own cost three. These exhaust $P\cup\{a_h\}$, so no deletion
satisfies~\eqref{eq:fallback}. Necessarily $a_h\in P$, and $c_h(P)=5$.

\paragraph{Only the pair owner can receive cake.}
For every singleton owner $j\ne h$, its item costs at most two to $h$.
Consequently
\[
 -c_h(P)+p_h\leq-5+C<-3<-2
 \leq-c_h(O_j)+p_j.
\]
Thus $h$ envies every other agent. Condition~\eqref{eq:fallback}
requires $p_j=0$ for each $j\ne h$, and completeness gives $p_h=C$.

\paragraph{An omitted designated item forces a contradiction.}
The pair $P$ omits at least one of $x,y,z$; denote an omitted item by
$z$, relabeling the three types if necessary. Both type-$Z$ agents
are singleton owners, since the pair contains its owner's designated
item. At most one of them holds $\{z\}$. Choose the other, say $k$.
Its own cost is two, whereas its cost for $P$ is three. Therefore
\[
 -c_k(O_k)+p_k=-2<-3+C=-c_k(P)+p_h.
\]
Agent $k$ envies the pair owner, whose cake has positive value.
Neither~\eqref{eq:EF} nor~\eqref{eq:fallback} can hold. This
contradiction proves nonexistence throughout $(1,2)$.

\section{The remaining cake totals}

For $0\leq C\leq1$, give $\{x,y\}$ and all the cake to $X_1$, give
$\{z\}$ to $Z_1$, and assign $d_1,d_2,d_3,d_4$ singly to
$X_2,Y_1,Y_2,Z_2$, respectively. Every singleton owner can remove its
own item to obtain item cost zero, so comparisons toward cake-free
bundles satisfy~\eqref{eq:EF} or~\eqref{eq:fallback}.
The pair owner can retain $x$, of cost one, after deleting $y$;
this also handles all of its comparisons toward cake-free bundles.
No singleton owner envies the paid pair: the three remaining $X/Y$
agents have own cost two and pair cost five, $Z_1$ has own cost one
and pair cost three, and $Z_2$ has own cost two and pair cost three.
The inequalities $2\leq5-C$, $1\leq3-C$, and $2\leq3-C$ hold for
$C\leq1$. Hence the allocation is strong EFM, including both endpoints.

For $C=2$, give $\{x,y\}$ and all cake to $Z_1$, give $\{z\}$ to
$Z_2$, and assign the ordinary chores singly to the four $X/Y$ agents.
This allocation is fully envy-free. Agent $Z_1$ has net cost one,
tying $\{z\}$ and preferring its package to each ordinary singleton.
Agent $Z_2$ has cost one and values the paid pair at net cost one.
Every $X/Y$ agent has cost two, values every singleton at cost two,
and values the paid pair at net cost three.
For $C>2$, use the same item allocation and payments
\[
 p_{Z_1}=2+\frac{C-2}{6},\qquad
 p_i=\frac{C-2}{6}\quad(i\ne Z_1).
\]
Adding the same amount to every payment preserves full envy-freeness.
Equation~\eqref{eq:payments} realizes these payments as a complete
cake partition. This completes the proof of Theorem~\ref{thm:main}.

\section{Why the strong definition matters}

Remark~2.4 of~\cite{ALMS} distinguishes a weaker condition that retains
the owner's cake value in the deletion comparison. Our counterexample
does not establish failure of that weaker condition. Indeed, for
$1<C<2$, use the last section's $C=2$ item allocation and give all cake
to $Z_1$. Only $Z_1$ envies $Z_2$: its net cost is $3-C>1$.
After deleting one of its own two items, the weak comparison has net
cost $2-C<1$, so it passes. Every other comparison is already envy-free.
The strong item-only comparison instead has cost two against cost one
and fails. Thus the negative answer applies precisely to the stronger
definition specified in D63-72.

\section{Scope, checking, and provenance}

Theorem~\ref{thm:main} gives a negative answer to the entire universal
existence question in D63-72 and in Section~5 of~\cite{ALMS}, under
Definition~2.3. A single admissible instance without an allocation
suffices to refute that assertion. The additional endpoint constructions
establish the exact budget set for this instance. There is no external
mathematical theorem on which the proof depends.

The definition, its weak variant, and the source question were checked
against the cited primary preprint. A bounded search of later primary
literature and current repository submissions found no matching complete
counterexample. This is not an exhaustive novelty or priority claim.

The mathematics and exposition were produced by OpenAI GPT-6 in Codex.
The exact deployed snapshot and reasoning configuration were not exposed.
Separate Codex agents using inherited settings checked every nonexistence
case, strict monotonicity, the endpoint allocations, and the source scope.
These are model-assisted checks; independent human proof checking and
Lean formalization have not been performed. The human supplied research
and publication instructions; submission is by \texttt{@jbaelaw}.

Final repository preparation followed Prompt A of the contribution policy,
then the final mathematical and compliance audit followed Prompt B.
The companion repository supplies the source, PDF, review records, and
an exhaustive structural checker. Its finite checks corroborate the hand
proof; they do not discretize cake or replace the proof over all cake
partitions. The initial research preceded this final preparation process.
\url{https://github.com/jbaelaw/strong-efm-nonadditive-counterexample}.

\section{Completion Estimate}
\noindent\textbf{COMPLETION: 100\%}
The explicit counterexample refutes the original universal strong-EFM
assertion. No mathematical case remains in that disproof. This completion
estimate records the manuscript's scope claim, not independent acceptance,
proof confidence, or a certification of novelty.

\begin{thebibliography}{9}
\bibitem{ALMS}
Haris Aziz, Xinhang Lu, Simon Mackenzie, and Mashbat Suzuki.
\emph{Fair Division with Indivisible Goods, Chores, and Cake}.
arXiv:2511.04891v1, 7 November 2025.
Definition~2.3, Remark~2.4, and Section~5.
\url{https://arxiv.org/abs/2511.04891v1}.
\end{thebibliography}
\end{document}
